% New material for the crossing-cores paper: a proved supermultiplicativity
% lemma with its corollaries, three exclusion lemmas, and two conjectures
% with exact evidence through order 22.  Drop-in sections; theorem
% environments as in the main file.

\section{Concatenation and supermultiplicativity}

For double-occurrence words $w_1$, $w_2$ on disjoint alphabets, write
$w_1w_2$ for the concatenation, read as a cyclic word.

\begin{lemma}[Profile factorization]\label{lem:profile}
For a pairing $\pi$ let
\[
f_\pi(x)\;=\;\operatorname{vol}_k\{t:\ 0\le x+S_\ell(t)\le 1,\ 0\le\ell<2k\}
\]
be the $x_0$-section volume of $\mathcal P_T(\pi)$, so that
$p_T(\pi)=\int_0^1 f_\pi$.  Then
\[
f_{w_1w_2}(x)\;=\;f_{w_1}(x)\,f_{w_2}(x),\qquad 0\le x\le1 .
\]
Moreover each $f_\pi$ is symmetric about $x=\tfrac12$ and log-concave.
\end{lemma}

\begin{proof}
Every pair contributes $+t_c$ and $-t_c$ to the walk, so after the last
slot of $w_1$ the walk has returned to $x_0$, and the slot forms of the
$w_2$-part of the concatenation involve only the $w_2$-variables.  The
constraint system therefore splits into the $w_1$-system and the
$w_2$-system, coupled only through the common $x_0$; Fubini in the two
independent blocks of $t$-variables gives the product formula.

Symmetry: the involution $(x_0,t)\mapsto(1-x_0,-t)$ maps every slot form
$W_\ell$ to $1-W_\ell$ and hence preserves $\mathcal P_T(\pi)$; it sends
the section at $x$ to the section at $1-x$.

Log-concavity: sections of a convex body by parallel hyperplanes have
$\frac1k$-concave, in particular log-concave, volumes, by
Brunn--Minkowski.
\end{proof}

\begin{lemma}[Supermultiplicativity]\label{lem:super}
For double-occurrence words on disjoint alphabets,
\[
p_T(w_1w_2)\;\ge\;p_T(w_1)\,p_T(w_2).
\]
The same holds for $p_H$ when $w_1$, $w_2$ and hence $w_1w_2$ are parity
mixing.
\end{lemma}

\begin{proof}
By Lemma~\ref{lem:profile}, $p_T(w_1w_2)=\int_0^1 f_{w_1}f_{w_2}$.  Both
factors are log-concave, hence unimodal, and symmetric about $\tfrac12$;
each is therefore a nondecreasing function of $-|x-\tfrac12|$.  Two
nondecreasing functions of the same variable are nonnegatively correlated
(Chebyshev's correlation inequality), so
$\int_0^1 f_{w_1}f_{w_2}\ge\int_0^1 f_{w_1}\int_0^1 f_{w_2}$.
The Hankel statement follows from Theorem~4.1, since concatenation of
parity-mixing words along even cut points is parity mixing.
\end{proof}

\begin{remark}
Equality in Chebyshev's inequality forces one profile to be constant, and
no nonempty word has constant profile (in every computed example
$f_\pi(0)<f_\pi(\tfrac12)$; at $x_0=0$ the section satisfies the extra
constraint $t_c\ge0$ for the first opened pair), so the inequality is
strict for nonempty words in all examples; we do not need the sharp
equality classification here.

The lemma explains the shape of the extremal hierarchy: the unique
four-pair Toeplitz maximizer $12123434$ is the concatenation of two copies
of the two-pair maximizer, with
$p_T=\tfrac9{20}\ge(\tfrac23)^2=\tfrac49$, and the six-pair parity-mixing
maximum $\tfrac{23}{90}$ is attained by the concatenation of two
three-pair cores of volume $\tfrac12$, with
$\tfrac{23}{90}\ge\tfrac14$.  Iterating the lemma on cube-cycle cores
gives persistent-volume families with an exact lower bound
$\prod_i 2/(r_i+1)$.
\end{remark}

\begin{lemma}[Cube-cycle profile]\label{lem:cube}
Let $w_r=12\cdots r\,12\cdots r$ be the fully crossing $r$-pair word.  Then
its section profile is
\[
f_{w_r}(x)\;=\;\int_0^1\min(x+y,\,2-x-y,\,1)^{\,r-1}\,dy
\;=\;\frac{1-x^r}{r}+\frac{1-(1-x)^r}{r},
\]
and consequently
\[
p_T(w_r)\;=\;\int_0^1 f_{w_r}\;=\;\frac{2}{r+1}.
\]
\end{lemma}

\begin{proof}
The slot forms of $w_r$ are the prefix sums $x_0+t_1+\dots+t_j$
($1\le j\le r$) and the suffix sums $x_0+t_{j+1}+\dots+t_r$
($1\le j\le r-1$).  Substitute $u_j=x+t_1+\dots+t_j$; the map
$t\mapsto u$ is unimodular, the prefix constraints become $u_j\in[0,1]$,
and the suffix constraints become $u_r+x-u_j\in[0,1]$ for $j<r$.  For
fixed $u_r=y$, each of $u_1,\dots,u_{r-1}$ ranges independently over
$[0,1]\cap[x+y-1,\,x+y]$, an interval of length
$\min(x+y,\,2-x-y,\,1)$.  Integrating over $y\in[0,1]$ gives the first
formula; splitting at $x+y=1$ gives the closed form.  For the volume,
with $s=x+y$ the sum of two independent uniforms,
\[
p_T(w_r)=\int_0^2\min(s,2-s,1)^{r-1}\,\bigl(1-|s-1|\bigr)\,ds
=2\int_0^1 s^{r-1}\cdot s\,ds=\frac{2}{r+1}. \qedhere
\]
\end{proof}

Combining Lemmas~\ref{lem:super} and~\ref{lem:cube}: for any concatenation
of cube cycles of sizes $r_1,\dots,r_m$,
\[
p_T\;\ge\;\prod_{i=1}^m\frac{2}{r_i+1},
\]
an explicit persistent-volume family with, by
Lemma~\ref{lem:profile}, exactly computable volumes
$\int_0^1\prod_i f_{w_{r_i}}$.

\begin{lemma}[Profile calculus]\label{lem:calculus}
Let $\pi$ be a pairing with section profile $f_\pi$ as in
Lemma~\ref{lem:profile}.
\begin{enumerate}
\item[(i)] (Ear neutrality.)  Inserting an ear into any gap of $\pi$
leaves the profile unchanged; hence $f_\pi=f_{\operatorname{core}(\pi)}$
pointwise, not only after integration.
\item[(ii)] (Flattening.)  If $w$ is a complete word and $a$ a new letter,
then $f_{awa}(x)=p_T(w)$ for every $x\in[0,1]$: nesting under a single
pair flattens the profile to the constant $p_T(w)$.
\item[(iii)] (Gap marginals.)  For a word $C$ and a gap $g$, let
$\rho_{C,g}$ be the pushforward to $[0,1]$ of Lebesgue measure on
$\mathcal P_T(C)$ under the slot form $W_g$.  Then $\rho_{C,g}$ has total
mass $p_T(C)$, is symmetric about $\tfrac12$ and log-concave, and for
every complete word $w$ on a disjoint alphabet inserted contiguously at
gap $g$,
\[
p_T\bigl(C[w@g]\bigr)\;=\;\int_0^1 f_w(u)\,\rho_{C,g}(du)\;\ge\;p_T(C)\,p_T(w).
\]
\end{enumerate}
\end{lemma}

\begin{proof}
(i) An ear inserted at a gap where the walk sits at $W_g$ contributes one
variable $t$ constrained only by $W_g+t\in[0,1]$, which is an interval of
length one for every admissible $W_g$; integrating it out gives the
factor $1$ at every fixed $x_0$.  This is the argument of Theorem~3.4
read at fixed $x_0$.

(ii) Opening $a$ moves the walk to $u=x+t_a$, which ranges over $[0,1]$
with unit Jacobian; the inner word runs from $u$ and returns to $u$, and
closing $a$ returns to $x\in[0,1]$.  Hence
$f_{awa}(x)=\int_0^1 f_w(u)\,du=p_T(w)$.

(iii) Conditioning on the variables of $C$, the inserted word starts and
ends at the position $W_g$ and contributes $f_w(W_g)$; integrating over
$\mathcal P_T(C)$ gives the formula.  The involution
$(x_0,t)\mapsto(1-x_0,-t)$ preserves $\mathcal P_T(C)$ and maps $W_g$ to
$1-W_g$, so $\rho_{C,g}$ is symmetric; it is log-concave as the marginal
of the indicator of a convex body (Pr\'ekopa).  Both $f_w$ and the
density of $\rho_{C,g}$ are then nondecreasing functions of
$-|u-\tfrac12|$, and Chebyshev's correlation inequality on $[0,1]$ gives
the bound.
\end{proof}

For cube cycles the gap marginals are not merely explicit; they are the
profile itself.

\begin{corollary}[Self-similarity of cycle gaps]\label{cor:selfsim}
For every gap $g$ of the fully crossing word $w_r$,
\[
\rho_{w_r,g}\;=\;f_{w_r},
\]
so for any complete word $w$ on a disjoint alphabet,
\[
p_T\bigl(w_r[w@g]\bigr)\;=\;\int_0^1 f_{w_r}f_w ,
\]
independently of which gap is used.
\end{corollary}

\begin{proof}
In the coordinates of Lemma~\ref{lem:cube} the slot form at an opener
gap is $u_j$, which for fixed $u_r=y$ ranges over an interval of length
$\min(x+y,2-x-y,1)$ together with the other $r-2$ free prefix
coordinates.  Writing $s=x+y$, whose density is $1-|s-1|$ on $[0,2]$,
the marginal is
$\rho(v)=\int_v^{v+1}\min(s,2-s,1)^{\,r-2}(1-|s-1|)\,ds$.  On $[0,2]$
one has $1-|s-1|=\min(s,2-s,1)$, so the integrand is
$\min(s,2-s,1)^{\,r-1}$ and $\rho=f_{w_r}$ by
Lemma~\ref{lem:cube}.  The rest is Lemma~\ref{lem:calculus}(iii).
\end{proof}

At $r=2$ this reads $f_{abab}(v)=\tfrac12+v(1-v)=\tfrac34-(v-\tfrac12)^2$,
so the four-pair extremal volume is an $L^2$ norm,
\[
\frac{9}{20}\;=\;\int_0^1 f_{abab}(v)^2\,dv ,
\]
consistent with $f_{12123434}=f_{abab}^{\,2}$ and $p=\int f$.  Likewise
$\int_0^1f_{w_3}f_{abab}=\tfrac{61}{180}$ and
$\int_0^1f_{w_4}f_{abab}=\tfrac{19}{70}$, matching the measured
insertion values at every gap of $abcabc$ and $abcdabcd$ respectively,
and $\int_0^1f_{w_3}^2=\tfrac{23}{90}$ is the six-pair parity-mixing
maximum.

The affine equivalence of Theorem~4.1 fixes the initial coordinate
$x_0$ and acts unimodularly on the fiber, so it identifies sections and
not only volumes: for every parity-mixing pairing the Hankel and
Toeplitz profiles are equal pointwise, $f^H_\pi(x)=f^T_\pi(x)$ for all
$x\in[0,1]$.  Every profile statement in this section therefore
transfers verbatim to the Hankel ensemble on the parity-mixing class.

Part (iii) contains Lemma~\ref{lem:super} (the seam gap) and shows that
the paper's inflation structure extends from ears to arbitrary
contiguous insertions, with an exact formula.  For cycle cores the gap
marginals coincide across gaps; for example every gap of $abab$ receives
$cdcd$ with volume $\tfrac9{20}$ and every gap of $abcabc$ receives
$dede$ with volume $\tfrac{61}{180}$, while the path-of-squares core
$abacbdcd$ has two gap classes, receiving $efef$ with volumes
$\tfrac14$ and $\tfrac{313}{1260}$.  Since ears are profile-neutral, the
Catalan inflation of Theorem~6.1 holds at the level of profiles:
$\sum_{\pi}f_\pi(x)=\mathrm{Cat}_k+\sum_{r}\binom{2k}{k-r}K_r(x)$ with
$K_r(x)=\sum_{\rho\in\mathcal R_{T,r}}f_\rho(x)$, and the generating
function identities of Theorem~6.4 hold for each fixed $x$.

\begin{lemma}[Supermultiplicativity of the primitive masses]\label{lem:kappa}
For all $r,s\ge2$,
\[
\kappa_{r+s}^{T}\;\ge\;\kappa_{r}^{T}\,\kappa_{s}^{T},
\qquad
\kappa_{r+s}^{H}\;\ge\;\kappa_{r}^{H}\,\kappa_{s}^{H}.
\]
In particular $\lim_r(\kappa_r^E)^{1/r}$ exists in $(0,\infty]$ for both
ensembles.
\end{lemma}

\begin{proof}
Fix cyclic positions $1,\dots,2(r+s)$.  For a reduced $r$-pair word $w_1$
on the first $2r$ positions and a reduced $s$-pair word $w_2$ on the rest,
on disjoint alphabets, the concatenation $w_1w_2$ is again reduced: linear
adjacencies inside each block are a subset of that block's cyclic
adjacencies, and the two seam adjacencies join letters of different
alphabets.  The map $(w_1,w_2)\mapsto w_1w_2$ is injective, since the
concatenation determines its restriction to the first $2r$ positions.
Therefore, using Lemma~\ref{lem:super} termwise,
\[
\kappa_{r+s}\;\ge\sum_{(w_1,w_2)}p(w_1w_2)\;\ge\;
\Bigl(\sum_{w_1}p(w_1)\Bigr)\Bigl(\sum_{w_2}p(w_2)\Bigr)
=\kappa_r\,\kappa_s .
\]
For the Hankel case both blocks have even length, so parity mixing is
preserved and Theorem~4.1 applies termwise.  Existence of the limit is
Fekete's lemma applied to $\log\kappa_r$.
\end{proof}

\begin{remark}\label{rmk:kappaasym}
The exact values are far above the bound
($\kappa_8=\tfrac{687080}{9}$ against $\kappa_4^2=\tfrac{31684}{225}$),
reflecting that only cut-separated concatenations are counted.  The data
suggest the sharper asymptotic
$\kappa_{r+1}/\bigl((r+1)\kappa_r\bigr)\to2\rho^2$ with $\rho$ the
Sen--Vir\'ag constant: the exact ratios $1.4833$, $1.4213$, $1.4002$,
$1.3899$, $1.3844$ have geometrically shrinking decrements and
extrapolate to $1.3734$, within $10^{-3}$ of $2\rho^2=1.37414$.  Non-separated interleavings satisfy no such inequality: $12341234$
interleaves $1212$ and $3434$ and has
$p_T=\tfrac25<\tfrac49=p_T(1212)\,p_T(3434)$ - simultaneously open pairs
couple the corridor constraints destructively.  So the binomial
enhancement of the lemma is false, and the finer $r!\gamma^r$ scaling
must come from the boundary-local structure of the Catalan transform
rather than from correlation inequalities.
\end{remark}

\begin{proposition}[Cluster equation]\label{prop:cluster}
Let $F(z;x)=\sum_k F_k(x)z^k$ with $F_k(x)=\sum_{\pi\in P_2(2k)}f_\pi(x)$
the profile sum over all pairings, read as linear words, and let
$J(z;x)$ be the same sum restricted to crossing-indecomposable words
(words that are neither a nontrivial concatenation of complete blocks
nor of the nested form $awa$).  Then, as formal power series,
\[
F(z;x)\;=\;\frac{1}{\,1-z\,M(z)-J(z;x)\,},
\qquad
M(z)\;=\;\int_0^1F(z;x)\,dx ,
\]
where $M$ is the even-moment generating function of $\gamma_T$.  With
$J\equiv0$ this is the Catalan equation $C=1/(1-zC)$, so $J$ is
precisely the crossing correction to the semicircle law.
\end{proposition}

\begin{proof}
Every nonempty linear word factors uniquely as a concatenation of
indecomposable complete words, and profiles multiply over concatenation
(Lemma~\ref{lem:profile}), so $F=1/(1-I)$ with $I$ the indecomposable
profile sum.  An indecomposable word is either of the form $awa$ with
$w$ complete - by the flattening rule its profile is the constant
$p_T(w)$, and summing over $w$ and the label of the enclosing pair gives
the term $z\,M(z)$ - or crossing-indecomposable, giving $J$.  Averaging
the first identity over $x\in[0,1]$ gives the second.
\end{proof}

The identity was verified exactly through order $z^4$ at $x=\tfrac12$:
series inversion of the computed $F_k(\tfrac12)=1,1,\tfrac{11}4,
\tfrac{47}4,\tfrac{12823}{192}$ yields
$J_k(\tfrac12)=0,\tfrac34,\tfrac{55}{12},\tfrac{6139}{192}$, matching
the direct enumeration of crossing-indecomposable words.  Iterating the
same decomposition inside the gaps of crossing-indecomposable words
expands $J$ over gap-irreducible atoms with gaps refilled by $F$
evaluated at the gap positions; the gap-marginal formula of
Lemma~\ref{lem:calculus}(iii) is the first moment of this expansion.
Note that section profiles are not invariant under rotation of the
cyclic word - rotating moves the basepoint - so all sums here are over
linear words; only the integrated volumes are rotation invariant.

\subsection*{The calculus closes on a class}

Lemmas~\ref{lem:profile}, \ref{lem:calculus}(ii) and~\ref{lem:cube}
together are an algorithm on the class of words generated from fully
crossing cycles by concatenation and by nesting under a single pair:
\[
W\;::=\;w_r\ \bigm|\ WW\ \bigm|\ aWa .
\]

\begin{proposition}[Closed evaluation]\label{prop:closure}
For every word in this class the section profile is a polynomial in $x$,
computed from
\[
f_{w_r}(x)=\frac{1-x^r}{r}+\frac{1-(1-x)^r}{r},\qquad
f_{UV}=f_Uf_V,\qquad
f_{aUa}\equiv\int_0^1 f_U ,
\]
and $p_T$ is the integral of that polynomial.  Every such volume is
therefore an explicit rational number, obtained without any polytope
computation.
\end{proposition}

The class is larger than its grammar suggests.  By
Corollary~\ref{cor:selfsim} the gap marginal of a cycle is its own
profile, so inserting a word $w$ into any gap of $w_r$ has volume
$\int_0^1f_{w_r}f_w$ -- which is exactly the concatenation formula.
Insertion into a cycle is therefore volume-equivalent to concatenation,
and every word obtained by hanging cycles off the gaps of other cycles
already has its volume realised inside the grammar.  This is why the
gap-independence observed for $abab$, $abcabc$ and $abcdabcd$ was not a
coincidence.

Reducedness needs one remark.  Read cyclically, $aWa$ places the two
endpoints of $a$ at positions $0$ and $2|W|+1$, which are adjacent, so a
word that is entirely a nesting is an ear and not a reduced core -- its
volume equals $p_T(W)$, as ear deletion requires.  Nesting contributes
reduced words only inside a concatenation, where the other factor
separates the endpoints; that is exactly how $\omega_l$ and $\eta_l$ are
built.

Because evaluation is now a polynomial integral, the maximum over the
whole class at a given number of pairs can be computed far past the
range of any polytope enumeration.  Doing so gives
\[
r=2:\ \tfrac23,\qquad r=3:\ \tfrac12,\qquad r=4:\ \tfrac9{20},\qquad
r=5,6,\dots,16:\ \tfrac49 ,
\]
the small cases reproducing the constants of Corollary~9.3 and
Theorem~9.1 and the rest agreeing with
Proposition~\ref{prop:family49} at every size up to sixteen pairs.  This
is a maximum over the grammar class rather than over all reduced cores,
and the class is thin: comparing against the complete reduced-core
sweeps, it realises $3$ of the $6$ distinct volumes at $r=5$, $6$ of
$19$ at $r=6$ and $8$ of $46$ at $r=7$, so its share of the spectrum
falls as $r$ grows.  Two things make the comparison worth making
anyway.  Every volume the class produces does occur in the true
spectrum -- $3/3$, $6/6$, $8/8$ at those sizes, so the calculus
introduces no spurious values -- and at each of $r=5,6,7$ the class
maximum equals the global maximum.  The evidence for $r\ge8$ is
therefore of the form: the extremum is known to lie in this class at
every size where the comparison can be made, and inside the class it
stays at $4/9$ through sixteen pairs.

What lies outside is also identifiable, and it says what the calculus is
missing.  The cores realising the volumes the class does not reach are
words such as $abacbdeced$, $abacbdcede$ and $abcadbecde$ at five pairs,
or $ababcdefefcd$ and $abacbdefcfed$ at six: in each the cycles of the
prefix-state cactus are glued at an interior shared state rather than at
the base point.  Concatenation joins blocks at the empty state and
nesting places one inside a single pair; neither produces a gluing at an
interior vertex, and for such a word the profile is not a product -- the
path-of-squares core $abacbdcd$ is the smallest example, and it is
exactly the case where the gap marginals differ from gap to gap
(volumes $\tfrac14$ and $\tfrac{313}{1260}$ for the same inserted word).
An evaluation rule for gluing two cycles at an interior state would
reach many of these; $\tfrac{313}{1260}$ duly appears among the six-pair
volumes outside the present class.  It would not reach all of them,
because the block data of Lemma 7.1 is not a complete invariant for the
volume.  The words
\[
abcdabcd,\qquad abacbdcd,\qquad abcadbcd
\]
all have eight distinct prefix states on four coordinates, so each has a
single cycle block of length eight using four labels -- the same data in
the classification -- yet their volumes are
\[
\tfrac25,\qquad \tfrac{11}{30},\qquad \tfrac{11}{30}.
\]
What separates them is how the cycle sits in the cube: the first walks
monotonically up and then down, the others zigzag.  So a complete
evaluation needs the embedded cycle, not merely its length and label
set.

The other obvious coarse invariant fails too.  Grouping all $105$
order-eight pairings by their height sequence -- the number of open
chords after each slot, that is the Dyck path -- gives $14$ classes, and
$13$ of them contain more than one volume; the class of
$(1,2,1,0,1,2,1,0)$ carries $1$, $\tfrac23$ and $\tfrac9{20}$ at once.
Note that $abacbdcd$ and $abcadbcd$ have different height sequences and
the same volume $\tfrac{11}{30}$, while words sharing a height sequence
routinely differ, so the Dyck path neither determines the volume nor is
determined by it.  The volume depends on which chord each down-step
closes, not on the height profile -- which is one reason these numbers
have resisted a formula, and why the profile calculus buys what it does
only on a structured subclass.

Among coarse statistics the crossing number does best.  At order eight
it takes seven values and five of the seven classes are pure; the
failures are exactly at two crossings, where $\tfrac12$ and $\tfrac9{20}$
both occur, and at three, where $\tfrac12$, $\tfrac{11}{30}$ and
$\tfrac25$ all occur.  The nesting number is nearly useless -- six of its
seven classes carry several volumes, and the class of nesting number
zero already carries all six values -- and refining to the pair
(crossings, nestings) still leaves seven impure classes out of
twenty-eight.  That the crossing number is so nearly a complete
invariant, and fails only in the middle of its range, invites the obvious
refinement: the chord-intersection graph $G_\times(\pi)$ of
Proposition~4.3 itself, rather than its edge count.  This is the
interlacement graph of the diagram, so the question is whether $p_T$
factors through the circle graph.

It very nearly does.  Grouping pairings by the isomorphism class of
$G_\times(\pi)$ gives, at order six, four classes and all pure; at order
eight, eleven classes and all pure, with $K_3\mapsto\tfrac12$ and
$K_4\mapsto\tfrac25$ in agreement with $p_T(w_r)=2/(r+1)$ for
$G_\times(w_r)=K_r$, and the empty graph $\mapsto1$.  At order ten there
are thirty-four classes and exactly one is impure: the graph on five
vertices with two disjoint edges and one isolated vertex carries both
$\tfrac9{20}$ and $\tfrac49$.  The witnesses are
\[
aabcbcdede\ \mapsto\ \tfrac9{20},\qquad
ababcdedec\ \mapsto\ \tfrac49 ,
\]
the second being $\omega_1$.  They have the same interlacement graph and
differ only in what the fifth chord does: in $\omega_1$ it nests the
second crossing pair, in the other it is an ear nesting nothing.  So the
volume is not a function of the circle graph alone.

The failure is not arbitrary, and separating the classes by connectivity
says where it lives.  Through order ten there are twenty-nine
isomorphism classes with $G_\times(\pi)$ connected --- two at order six,
six at order eight, twenty-one at order ten --- and every one of them is
pure.  Of the twenty disconnected classes in the same range, nineteen
are pure and the twentieth is the one displayed above, which is the only
impurity anywhere through order ten.  This is what the witnesses
suggest: when $G_\times(\pi)$ is
disconnected the diagram may realise its components either side by side
or with one nested inside another, and the volume tells those apart,
whereas a connected graph leaves no such freedom.  We record the
resulting question.

\begin{conjecture}
If $G_\times(\pi)$ is connected, then $p_T(\pi)$ depends only on the
isomorphism class of $G_\times(\pi)$.
\end{conjecture}

\begin{remark}
Part of this is classical.  By Bouchet's theorem a circle graph that is
prime for split decomposition has an essentially unique chord diagram,
so for prime $G_\times(\pi)$ the conjecture is immediate --- the graph
determines the diagram, never mind the volume.  Disconnection is the
crudest way for primeness to fail, and it is the way our witness fails.
What the order-ten data leaves open, and what would give the conjecture
real content, is the connected but non-prime range: we have not checked
which of the twenty-one connected classes are prime, and it is exactly
the imprime ones that would make this more than a restatement.
\end{remark}

The ingredient such a rule needs is a two-point kernel: the section
volume with both the base point and the shared state held fixed.  For
the square it is elementary.  The walk visits $x$, $u_1$, $u_2$ and
$u_2-u_1+x$, so fixing $x$ and $u=u_1$ and integrating out $u_2$ gives
\[
K_{abab}(x,u)\;=\;\bigl(1-|u-x|\bigr)_+ ,
\]
and $\int_0^1K(x,u)\,du=f_{abab}(x)$, $\iint K=\tfrac23$, both checked
exactly.  Profiles are the one-point marginals of these kernels;
concatenation and nesting only ever need the one-point data, which is
why they close, while gluing at an interior state needs $K$ itself.

The class also supplies a check of the calculus against the lattice
geometry.  Corollaries~5.2 and~5.3 say that the normalised volume
$(r+1)!\,p_T$ of an $r$-pair core is an integer and, by the parity
refinement, an even one.  Those statements come from Gorenstein
reciprocity and the mod-two $h^*$-identity; the grammar volumes come
from integrating polynomials and never touch a polytope.  All $628$
volumes the class produces for $2\le r\le12$ have even integral
normalised volume, with no exception.

The extremum is isolated, and by a fixed amount.  The three largest
volumes in the class are $\tfrac49$, $\tfrac13$, $\tfrac3{10}$ for every
$r\ge7$ (at $r=5,6$ the runners-up are $\tfrac{61}{180}$ and
$\tfrac13$), so the gap below the maximum is $\tfrac49-\tfrac13=\tfrac19$
independently of $r$.  This is a different phenomenon from the quantized
improvement of Corollary~9.2, whose gap shrinks like $2/(r+1)!$: within
this class the top of the volume spectrum does not accumulate at $4/9$
but sits a fixed distance below it.

The same computation runs on the Hankel side.  Since $w_r$ is parity
mixing exactly for odd $r$, and both concatenation (an even shift) and
nesting (which flips every parity) preserve parity mixing, the
parity-mixing subclass is the same grammar with odd cycles only.  Its
maxima are
\[
r=3:\ \tfrac12,\qquad r=5:\ \tfrac13,\qquad r=6:\ \tfrac{23}{90},\qquad
r=7,8,\dots,18:\ \tfrac14 ,
\]
reproducing Corollary~9.3 at three pairs, Proposition~10.2 at five, the
six-pair maximum, and then the $\eta_l$ value at every size up to
eighteen.  One gap remains: the grammar contains no four-pair
parity-mixing word, so it misses the extremal $12314324$ of volume
$11/30$; the class is rich enough to carry the asymptotic extremals but
not every small one.

The class already contains the extremal families: $abab\,cdcd$ gives
$\tfrac9{20}$, the two-square chains $\omega_l$ give $\tfrac49$, the
$w_3$ chains $\eta_l$ give $\tfrac14$, and $w_3w_3$ gives
$\tfrac{23}{90}$.  It also produces values that would be laborious
otherwise, all confirmed against the exact engine: $abcabc\,dede$ gives
$\tfrac{61}{180}$, $abab\,cdcd\,efgefg$ gives $\tfrac{13}{56}$,
$abcabc\,dede\,fghfgh$ gives $\tfrac{166}{945}$, and
$abcdbcda\,efghefgh$ gives $\tfrac15$.  In fourteen randomly generated
words of the grammar the closed evaluation and the polytope volume agreed
exactly in every case.

\section{Exclusion lemmas for the limiting laws}

\begin{lemma}\label{lem:notgauss}
Neither $\gamma_T$ nor $\gamma_H$ is a variance mixture of centered
Gaussian laws.
\end{lemma}

\begin{proof}
If $X=\sqrt V Z$ with $Z$ standard Gaussian and $V\ge0$ independent, then
$m_4=3\,\mathbb E[V^2]\ge3(\mathbb E V)^2=3m_2^2=3$.  But
$m_4(\gamma_T)=\tfrac83<3$ and $m_4(\gamma_H)=2<3$.
\end{proof}

\begin{lemma}\label{lem:hankelnotsemi}
$\gamma_H$ is not a scale mixture of semicircle laws.
\end{lemma}

\begin{proof}
If $\gamma_H=\operatorname{law}(\sqrt V S)$ with $S$ standard semicircular,
then $\mathbb E[V^k]=m_{2k}(\gamma_H)/\mathrm{Cat}_k$.  Here
$\mathbb E[V]=1$ and $\mathbb E[V^2]=m_4/2=1$, so $V$ is degenerate at $1$
and $\gamma_H$ would be the semicircle law, contradicting
$m_6(\gamma_H)=\tfrac{11}2\ne5$.
\end{proof}

\begin{lemma}\label{lem:notq}
Neither law is a $q$-Gaussian.  For $\gamma_T$, $m_4=2+q$ forces
$q=\tfrac23$, and then $m_6=5+6q+3q^2+q^3=\tfrac{287}{27}\ne11$; for
$\gamma_H$, $m_4=2$ forces $q=0$, and $m_6=\tfrac{11}2\ne5$.
\end{lemma}

Lemma~\ref{lem:notgauss} rules out any exact conditional-Gaussian or
symbol-averaging representation of either law.
Lemma~\ref{lem:hankelnotsemi} is consistent with the non-unimodality of
$\gamma_H$: a scale mixture of semicircle laws is unimodal.

\section{The global $4/9$ bound: extremals and exhaustive evidence}

Beyond the four-pair layer, Theorem~9.1's constant $9/20$ is attained only
at $k=4$, and the natural question is the sharp constant for $r\ge5$.  The
profile calculus supplies the extremal family in closed form.

\begin{proposition}[The $4/9$ family]\label{prop:family49}
For $l\ge1$ let
\[
\omega_l\;=\;abab\,c_1c_2\cdots c_l\,dede\,c_l\cdots c_2c_1 ,
\]
a reduced core with $r=l+4$ pairs (two square blocks joined by a chain of
$l$ nested pairs).  Then, for every $x\in[0,1]$,
\[
f_{\omega_l}(x)\;=\;\tfrac23\,f_{abab}(x),
\qquad\text{hence}\qquad
p_T(\omega_l)\;=\;\Bigl(\tfrac23\Bigr)^2=\frac49 .
\]
\end{proposition}

\begin{proof}
The block $c_1\cdots c_l\,dede\,c_l\cdots c_1$ is a complete word nested
under $c_1$; applying the flattening rule of
Lemma~\ref{lem:calculus}(ii) $l$ times reduces it to the constant profile
$p_T(dede)=\tfrac23$ (the two-pair crossing volume).  The whole word is
the concatenation of $abab$ with that block, so profiles multiply by
Lemma~\ref{lem:profile}, giving
$f_{\omega_l}=\tfrac23f_{abab}$ and
$p_T(\omega_l)=\tfrac23\int_0^1f_{abab}=\tfrac23\cdot\tfrac23$.
\end{proof}

\begin{proposition}[Exhaustive verification for $r\le7$]\label{prop:49exh}
Over all reduced Toeplitz cores with $r=5,6,7$ pairs -- that is $293$,
$3326$ and $44189$ labelled words, $29$, $196$ and $1788$ dihedral orbits
-- the maximum volume is exactly $4/9$, and it is attained on a single
dihedral orbit at each $r$, namely the orbit of $\omega_{r-4}$.
\end{proposition}

Together these give the sharp form of the conjecture: for every $r\ge5$,
\[
p_T(\rho)\;\le\;\frac49 ,
\]
with equality exactly on the orbit of $\omega_{r-4}$ -- an equality class
that is nonempty for every $r\ge5$, so the constant cannot be improved at
any size.

The general-$r$ statement reduces to a single finite slab lemma.  Recall
the proof of Theorem~9.1: reducedness forces $e\ge2$; for $e=2$
Corollary~7.2 gives $9/20$ (only at $k=4$) or $4/9$; and for $e\ge3$ the
paper bounds by $9/20$ or $2/5$.  Hence the global bound for $r\ge5$
follows at once from the following strengthening of Lemma~8.2, which
removes its criticality hypothesis.

\begin{conjecture}[Three-slab bound]\label{conj:threeslab}
Let $L=X_1+X_2-X_3$ and let $M\ne N$ be slot forms, neither a coordinate
nor $\pm L$.  Then
\[
\Pr\bigl\{|L|\le\tfrac12,\ |M|\le\tfrac12,\ |N|\le\tfrac12\bigr\}
\;\le\;\frac25 .
\]
\end{conjecture}

By Lemma~2.5 every slot form has coefficients in $\{0,\pm1\}$ with odd
support of size at least three, so for each ambient dimension $d$ this is
a finite exact computation of the same kind as Appendix~A and
Theorem~10.4.  For forms of minimal support the finiteness is absolute
and the conjecture is a theorem.

\begin{proposition}[Three triangles]\label{prop:threetri}
Let $M_1,M_2,M_3$ be pairwise distinct slot forms of support three, none
the negative of another.  Then
\[
\Pr\bigl\{|M_i(X)|\le\tfrac12,\ i=1,2,3\bigr\}\;\le\;\frac25 ,
\]
with equality exactly for the star system of Lemma~\ref{lem:star}.
\end{proposition}

\begin{proof}
A support-three form with coefficient sum one has two entries $+1$ and
one entry $-1$, so the configuration is a $3\times d$ matrix over
$\{0,\pm1\}$ whose rows each carry two $+1$ and one $-1$, with no zero
column.  Permuting coordinates does not change the probability, so a
configuration is determined by the multiset of its columns, each a
nonzero vector in $\{-1,0,1\}^3$, and the row conditions say that for
each row the multiplicities of the columns carrying $+1$ there sum to
two and those carrying $-1$ sum to one.  That integer system has exactly
$463$ solutions, of which $34$ are degenerate (two forms equal or
opposite).  Computing the remaining $429$ probabilities exactly gives
maximum $2/5$, attained by five configurations, all equivalent to
\[
M_1=U+X_1,\quad M_2=U+X_2,\quad M_3=U+X_3,\qquad U=X_5-X_4 ,
\]
the star of Lemma~\ref{lem:star} at $j=3$.  The next values attained are
$3/8$, $11/30$, $17/48$, $61/180$ and $1/3$.
\end{proof}

Since the star is itself a triangle system,
Proposition~\ref{prop:threetri} settles the conjecture on the set where
the extremum lives, and what remains is to rule out larger supports.
Exhaustive scans over forms of every admissible support give maximum
$3/8$ over all $55$ pairs at $d=4$, $2/5$ over all $741$ at $d=5$, and
$2/5$ over all $7021$ at $d=6$; the value $2/5$ is Lemma~8.2's
critical-triple constant.  Since $\tfrac25<\tfrac49$, the conjecture
yields
\[
p_T(\rho)\le\frac49\quad(r\ge5),\qquad
\text{equality exactly on the orbit of }\omega_{r-4},
\]
a complete proof of the global bound modulo one finite certificate.

\subsection*{The Hankel analogue: a sharp $1/4$}

The parity restriction forbids square blocks (Corollary~7.3), so the
smallest cycle available in the parity-mixing class is $w_3=abcabc$ with
$p_H(w_3)=\tfrac12$.  Replacing squares by $w_3$ in
Proposition~\ref{prop:family49} gives the Hankel analogue.

\begin{proposition}[The $1/4$ family]\label{prop:family14}
For $l\ge1$ let
\[
\eta_l\;=\;abcabc\,c_1\cdots c_l\,defdef\,c_l\cdots c_1 ,
\]
a reduced parity-mixing core with $r=l+6$ pairs.  Then
$f_{\eta_l}=\tfrac12 f_{w_3}$ and
\[
p_H(\eta_l)\;=\;p_T(\eta_l)\;=\;\Bigl(\tfrac12\Bigr)^2=\frac14 .
\]
\end{proposition}

\begin{proof}
Flattening (Lemma~\ref{lem:calculus}(ii)) applied $l$ times reduces the
nested block to the constant $p_T(w_3)=\tfrac12$; concatenation
multiplies profiles; and the Hankel value equals the Toeplitz value by
Theorem~4.1, the words being parity mixing.
\end{proof}

Computing the complete reduced parity-mixing layers gives the exact
maxima
\[
r=5:\ \tfrac13,\qquad r=6:\ \tfrac{23}{90},\qquad r=7:\ \tfrac14,
\qquad r=8:\ \tfrac14 ,
\]
the value at $r=8$ being new here (the whole eight-pair layer is $4738$
words in $198$ orbits, $92$ distinct volumes).  The equality classes are
also explicit: at $r=5$ the cube cycle $w_5$; at $r=6$ the single orbit
of $w_3w_3$, whose volume $\int_0^1f_{w_3}^2=\tfrac{23}{90}$ is again a
concatenation value; at $r=7$ exactly two orbits, the cube cycle $w_7$
and $\eta_1$; and at $r=8$ the single orbit of $\eta_2$ (the cube cycle
$w_8$ is not parity mixing -- $w_r$ is parity mixing precisely for odd
$r$).  Since $\eta_l$ exists for every $r\ge7$, this suggests the sharp
Hankel companion of Proposition~\ref{prop:family49}:
\[
p_H(\rho)\;\le\;\frac14\qquad(r\ge7),
\]
sharpening Corollary~10.5's $1/3$, which is attained only at $r=5$.  The
two bounds have the same shape: $4/9=p_T(w_2)^2$ on the Toeplitz side and
$1/4=p_H(w_3)^2$ on the Hankel side, each the square of the smallest
admissible cycle's volume.

\subsection*{One family behind every slab constant}

The separate constants appearing across Sections~8--10 are all values of
the cube-cycle family.  The fully crossing $r$-pair word $w_r$ has
exactly $r-1$ exceptional prefix-state forms, so writing $j=r-1$ for the
number of slabs, Lemma~\ref{lem:cube} gives
\[
p_T(w_{j+1})\;=\;\frac{2}{j+2}.
\]
Comparing with the constants in the paper:

\begin{center}
\begin{tabular}{c|c|l}
$j$ & $2/(j+2)$ & where the constant appears \\\hline
$1$ & $2/3$ & concentration of a three-uniform sum on a unit interval\\
$2$ & $1/2$ & Lemma 8.1, critical companions (equality)\\
$3$ & $2/5$ & Lemma 8.2, three critical slabs\\
$4$ & $1/3$ & Theorem 10.4, the four-slab theorem\\
$5$ & $2/7$ & the rank-five bound, open\\
$6$ & $1/4$ & predicted
\end{tabular}
\end{center}

Every constant matches, and by Lemma~\ref{lem:cube} the cube cycle
attains it at each level.  This suggests the single statement behind the
whole hierarchy.

\begin{lemma}[The star system attains $2/(j+2)$]\label{lem:star}
Let $U=X_a-X_b$ for two coordinates $a\ne b$, and let
$M_i=U+X_{c_i}$, $i=1,\dots,j$, for distinct coordinates $c_i$ outside
$\{a,b\}$ (signs immaterial).  Then
\[
\Pr\bigl\{|M_i|\le\tfrac12,\ 1\le i\le j\bigr\}
=\int_{-1}^{1}(1-|u|)^{j+1}\,du=\frac{2}{j+2}.
\]
This is the exceptional system of the fully crossing word $w_{j+1}$,
whose forms are $\Phi_0+\Phi_{j+1}-\Phi_i$, $1\le i\le j$.
\end{lemma}

\begin{proof}
$U$ has density $(1-|u|)_+$.  Given $U=u$, the constraint $|u+X_{c_i}|\le
\tfrac12$ puts $X_{c_i}$ in an interval of length $(1-|u|)_+$, and the
$X_{c_i}$ are independent, so the conditional probability is
$(1-|u|)^j$.  Integrating against the density of $U$ gives
$\int_{-1}^1(1-|u|)^{j+1}du=2/(j+2)$.
\end{proof}

Lemma~\ref{lem:star} recovers every constant in the table by a single
computation, and matches $p_T(w_{j+1})=2/(j+2)$ from
Lemma~\ref{lem:cube}, as it must.  The exhaustive scan at $j=3$ finds
exactly this configuration at the maximum: the nine attaining pairs at
$d=5$ all have the form $L=X_0+X_1-X_2$, $M=-U+X_4$, $N=-U+X_3$ with
$U=X_0-X_2$.

\begin{conjecture}[Slab hierarchy]\label{conj:slabhier}
For $j\ge1$ distinct exceptional slot forms $M_1,\dots,M_j$,
\[
\Pr\bigl\{|M_i(X)|\le\tfrac12,\ 1\le i\le j\bigr\}\;\le\;\frac{2}{j+2},
\]
with equality exactly for the exceptional system of the fully crossing
word $w_{j+1}$.
\end{conjecture}

The cases $j=1,2,4$ are theorems in the paper; Conjecture~\ref{conj:threeslab}
is the case $j=3$, and $j=5$ is the open rank-five bound.  Two reductions
make each case finite.  First, the coefficients are automatically in
$\{0,\pm1\}$ with odd support (Lemma~2.5).  Second:

\begin{lemma}[Private pairs may be deleted]\label{lem:anderson}
Let $M_1,\dots,M_j$ be slot forms and suppose two coordinates $c\ne d$
occur only in $M_1$, with opposite signs there.  Let $M_1'$ be $M_1$ with
those two entries deleted.  Then
\[
\Pr\bigl\{|M_i|\le\tfrac12\ \forall i\bigr\}\;\le\;
\Pr\bigl\{|M_1'|\le\tfrac12,\ |M_i|\le\tfrac12\ (i\ge2)\bigr\},
\]
and $M_1'$ is again a slot form (its support drops by two, staying odd,
and its coefficient sum is unchanged).
\end{lemma}

\begin{proof}
Write $Z=(M_1',M_2,\dots,M_j)$ and $W=\pm(X_c-X_d)$, so that the vector
of forms is $Z+(W,0,\dots,0)$ with $W$ independent of $Z$.  The event is
$Z+(W,0,\dots,0)\in K$ for the symmetric convex set
$K=[-\tfrac12,\tfrac12]^j$.  Now $Z$ is a linear image of a uniform
distribution on a cube, hence has a symmetric log-concave law, so
Anderson's inequality gives $\Pr\{Z+v\in K\}\le\Pr\{Z\in K\}$ for every
fixed $v$; averaging over $v=(W,0,\dots,0)$ gives the claim.
\end{proof}

The hypothesis of opposite signs is needed: deleting two private
coordinates of the same sign would change the coefficient sum, so the
reduced row is no longer a slot form and the induction does not close.
What Lemma~\ref{lem:anderson} does give is that an extremal
configuration has no form carrying two private coordinates of opposite
sign.  A form may still carry several private coordinates provided they
all have the same sign -- a support-three form, for instance, can have
its two $+1$ entries private -- and such configurations are not reduced
by the lemma.  Combined with the exhaustive scans below,
which settle every configuration living in dimension at most six, the
cases left open are those in dimension at least seven in which each
form's private coordinates all carry the same sign.  Exhaustive checks: at $j=3$ the maximum is $3/8$ over all $55$
admissible pairs at $d=4$ and $2/5$ over all $741$ at $d=5$, and the
deletion step never increased the probability in $55$ of $55$ tested
systems.  At $j=5$ the star system was confirmed to give exactly $2/7$,
and a search over five-subsets of admissible forms found nothing
exceeding it: exhaustively at $d=5$ (maximum $29/120$) and by random
sampling at $d=6$ and $d=7$ (maxima $37/160$ and $1891/10080$).

A sharper test than sampling is to ask whether the star can be improved
locally.  Starting from the star and replacing one of its forms by every
other admissible slot form, computing each result exactly, gives at
$j=3$ a best value of $2/5$ over $111$ replacements in $d=5$ and over
$351$ in $d=6$, and at $j=4$ a best value of $1/3$ over $464$
replacements in $d=6$ and $1384$ in $d=7$ -- $2310$ exact
recomputations in all, and in none of them does anything beat the star,
which is therefore an exact local maximum under single-form moves, also
when a spare coordinate outside its support is available.  The $j=4$ run
is also a check on the machinery rather than on the conjecture, since
Theorem 10.4 proves the bound $1/3$ there; a replacement exceeding it
would have contradicted a proved theorem, and none does.

Those sampling figures should be read with care, and their own numbers
say why.  Simulating random $j$-tuples for $j=4,5,6$ in dimensions up to
$j+4$ gives best observed values $0.301$, $0.208$ and $0.117$ against
bounds $1/3$, $2/7$ and $1/4$: the sampled maxima fall further below the
bound as $j$ and $d$ grow, so a random tuple is almost never
extremal-shaped and "no violation found" is weak evidence at large $j$.
The weight is carried by the exhaustive scans in low dimension and by
Proposition~\ref{prop:threetri}, which is a complete enumeration.  What
the simulations do confirm sharply is the equality case: the star system
reproduces $2/(j+2)$ to within Monte-Carlo error at every $j$ from $2$
to $7$, including the values $1/4$ and $2/9$ beyond the range where the
exact volumes were computed.

A route to the open case $j=5$ already exists inside the paper's own
apparatus.  Theorem~10.4 proves $j=4$ by writing $D=[-M\,|\,I_4]$,
observing that its column matroid is regular of rank four, and invoking
the rank-four classification ($M(K_5)$ restrictions plus
$M^*(K_{3,3})$).  The identical argument at $j=5$ needs the rank-five
regular matroid data, and the cosystole bounds of Douthitt, Israel and
Kennard cited there cover regular matroids of rank at most six.  So the
rank-five $2/7$ bound looks accessible by the paper's existing method
rather than a new one.

\section{Unimodality: a moment certificate}

Khinchin's theorem states that a symmetric law $\nu$ is unimodal about
the origin if and only if $\nu$ is the law of $UZ$ with $U$ uniform on
$[-1,1]$ independent of $Z$; equivalently, if and only if
\[
u_k\;:=\;(2k+1)\,m_{2k}(\nu),\qquad k\ge0,
\]
is the moment sequence of a positive random variable ($Z^2$).  This turns
unimodality into an exactly checkable Hankel positivity condition.

\begin{proposition}[Four-moment certificate for the Hankel law]\label{prop:hnonuni}
$\gamma_H$ is not unimodal.  Indeed, with
$u_0,\dots,u_4=1,\,3,\,10,\,\tfrac{77}2,\,\tfrac{843}5$,
\[
\det\begin{pmatrix}u_0&u_1&u_2\\u_1&u_2&u_3\\u_2&u_3&u_4\end{pmatrix}
=\frac{815}{4}-\frac{1812}{5}+155=-\frac{73}{20}\;<\;0 .
\]
\end{proposition}

\begin{proof}
If $\gamma_H$ were unimodal, Khinchin's theorem would give
$\gamma_H=\operatorname{law}(UZ)$ and hence $u_k=\mathbb E[(Z^2)^k]$, so
the Hankel matrix $(u_{i+j})_{i,j\le2}$ would be positive semidefinite.
The displayed determinant, computed from $m_2=1$, $m_4=2$, $m_6=\tfrac{11}2$
and $m_8=\tfrac{281}{15}$, is negative.
\end{proof}

Non-unimodality of $\gamma_H$ is known~\cite[Prop.~A.2]{BDJ}; the point
here is that the exact moments make it a four-line rational computation,
and that the same test is informative for the Toeplitz law, where
unimodality is open.  Two controls show the criterion has teeth: for the
semicircle law all computed minors are positive, and for the symmetric
two-point law $\tfrac12(\delta_{-1}+\delta_1)$ the criterion fails
already at $D_1=-4$.

\begin{proposition}[Toeplitz unimodality is not obstructed through order 22]
\label{prop:tuni}
For $\gamma_T$ the sequence $u_k=(2k+1)m_{2k}(\gamma_T)$ begins
$1,\,3,\,\tfrac{40}3,\,77,\,\tfrac{2724}5,\,4565,\dots$, and all six
leading minors of $(u_{i+j})$ and all six of $(u_{i+j+1})$ are strictly
positive.  Hence no obstruction to the unimodality of $\gamma_T$ is
visible in the moments through order $22$.
\end{proposition}

This is consistent with, and weaker than,
Conjecture~\ref{conj:mix}: a scale mixture of semicircle laws is
unimodal, so the semicircular-mixture conjecture implies unimodality of
$\gamma_T$, and the two mixture tests indeed pass together at every
computed order.

The status of the two laws is asymmetric here.  For $\gamma_H$,
non-unimodality is part of the original theorem of Bryc, Dembo and
Jiang~\cite{BDJ}, who use the convexity form of the definition -- a
symmetric law $\nu$ is unimodal when $x\mapsto\nu((-\infty,x])$ is
convex on $(-\infty,0)$ -- which is the form Khinchin's theorem
characterises, so
Proposition~\ref{prop:hnonuni} is a certificate for exactly that
statement.  For $\gamma_T$, Sen and Vir\'ag~\cite{SenVirag} proved
absolute continuity with density bounded by $8$, answering part of the
question raised in~\cite{BDJ}; unimodality is not settled there and we
are not aware of a later resolution.  Proposition~\ref{prop:tuni} says
the moments through order twenty-two contain no obstruction to it.

\section{Two conjectures}

\begin{conjecture}[Semicircular mixture]\label{conj:mix}
$\gamma_T$ is a scale mixture of semicircle laws: there is a probability
law $\mu$ on $(0,\infty)$ with
\[
\gamma_T=\operatorname{law}\bigl(\sqrt V\,S\bigr),\qquad V\sim\mu,\quad
\mathbb E[V^k]=\frac{m_{2k}(\gamma_T)}{\mathrm{Cat}_k}.
\]
\end{conjecture}

The mixing moments through order $22$ are
\begin{gather*}
1,\quad \frac43,\quad \frac{11}5,\quad \frac{454}{105},\quad
\frac{415}{42},\quad \frac{5960}{231},\quad \frac{108573}{1430},\\[2pt]
\frac{56067533}{225225},\quad \frac{3065263961}{3403400},\quad
\frac{17293570111}{4849845},\quad \frac{298983373961}{19492200},
\end{gather*}
and both Stieltjes--Hankel matrices built from them are positive definite:
all six leading minors of $(s_{i+j})$ and all six of $(s_{i+j+1})$ are
strictly positive.  A finitely supported mixing law would annihilate a
minor, so the evidence is nondegenerate at every available order.  In
generating-function form the claim reads $M_T(z)=\mathbb E[C(Vz)]$ with
$C$ the Catalan generating function, so Theorem~6.4 expresses the
primitive masses $\kappa_r^T$ as explicit functionals of $\mu$.  The
conjecture implies $\gamma_T$ is unimodal.  The mixing law is close to,
but is not, a Gamma law: $\Gamma(3,\tfrac13)$ matches $\mathbb E V$ and
$\mathbb E V^2$ and gives $\mathbb E V^3=\tfrac{20}9\ne\tfrac{11}5$.

\begin{conjecture}[Recurrence slope]\label{conj:slope}
Let $b_k^2$ be the Jacobi coefficients of $\gamma_T$, so that its
orthogonal polynomials satisfy $p_{k+1}=xp_k-b_k^2p_{k-1}$.  Then
\[
b_{k+1}^2-b_k^2\;\longrightarrow\;\rho^2\;=\;0.6870697\ldots,
\]
where $\rho=0.8288582\ldots$ is the Sen--Vir\'ag operator-norm constant.
\end{conjecture}

The exact values give $b_k^2=1,\ \tfrac53,\ \tfrac73$ (increments exactly
$\tfrac23$), then increments
$0.68952,\ 0.69342,\ 0.69200,\ 0.69103,\ 0.69046,\ 0.68979,\ 0.68939,\
0.68905$ through order $22$: after the peak at $k=4$ the sequence
decreases monotonically toward $\rho^2$.  The conjectured mechanism is
that $b_k^2\sim ck$ characterizes a Gaussian-type tail
$\exp(-x^2/2c)$ for the limiting law, and the extreme eigenvalues of the
normalized matrix then grow like $\sqrt{2c\log n}$, which is the
Sen--Vir\'ag asymptotic $\rho\sqrt{2\log n}$ precisely when $c=\rho^2$.
The conjecture is thus equivalent to the statement that the Gaussian tail
variance of $\gamma_T$ equals $\rho^2$.

The two conjectures are welded together by the following elementary but
rigorous implication.

\begin{proposition}[Increments control the mean volume]\label{lem:dyck}
Let $\nu$ be a symmetric law with finite moments and Jacobi coefficients
$b_j^2>0$ satisfying $b_{j+1}^2-b_j^2\to c>0$.  Then
\[
\Bigl(\frac{m_{2k}(\nu)}{(2k-1)!!}\Bigr)^{1/k}\;\longrightarrow\;c ,
\]
so Conjecture~\ref{conj:slope} forces the mean pairing-polytope volume
$\mathbb E_\pi[p_T(\pi)]$ to decay at the exact geometric rate $\rho^2$.
\end{proposition}

\begin{proof}
The lattice-path expansion gives
$m_{2k}=\sum_{P}\prod_i b_{h_i}^2$ over Dyck paths of semilength $k$,
where $h_1,\dots,h_k$ are the heights from which the down-steps fall.
For $b_j^2=j$ the sum is exactly $(2k-1)!!$: a down-step from height $h$
chooses which of $h$ open chords to close, so weighted paths biject with
matchings of $\{1,\dots,2k\}$.  Normalizing,
\[
\frac{m_{2k}}{(2k-1)!!}\;=\;\mathbb E\Bigl[\prod_{i=1}^{k}
\frac{b^2_{h_i}}{h_i}\Bigr],
\]
the expectation over the open-chord process $(h_j)_{0\le j\le 2k}$ of a
uniform random matching revealed left to right.  This process is an
explicit inhomogeneous Markov chain: after $j$ reveals at height $h$,
the next position closes an open chord with probability $h/(2k-j)$ (its
partner is exchangeable among the $2k-j$ unrevealed positions) and opens
a new chord otherwise.

Fix $\varepsilon\in(0,c)$ and choose $H$ with
$|b_h^2/h-c|\le\varepsilon$ for $h>H$; put
$B=\max_{h\le H}\max(b_h^2/h,\,h/b_h^2)\ge1$ and let $L$ be the number
of down-steps from heights $\le H$.  Then
\[
(c-\varepsilon)^{k}\,\mathbb E\bigl[\bigl(\tfrac{1}{B(c-\varepsilon)}\bigr)^{L}\bigr]
\;\le\;\frac{m_{2k}}{(2k-1)!!}\;\le\;
(c+\varepsilon)^{k}\,\mathbb E\bigl[\bigl(\tfrac{B}{c+\varepsilon}\vee1\bigr)^{L}\bigr].
\]
It remains to show $\mathbb E[x^{L}]=k^{O_x(1)}$ for every fixed
$x\ge1$, and $\mathbb E[y^{L}]\ge k^{-O_y(1)}$ for $y\in(0,1]$; both
follow from: $L\le \Lambda:=\#\{j:h_j\le H\}$, and $\Lambda$ is
$O_{H}(\log k)$ with tails controlled at every exponential scale.  Fix
$x\ge1$ and set $M=4xH$.  For $j\le 2k-M$ and $h\le H$ the chain steps
down with probability $h/(2k-j)\le H/M=1/(4x)$, so during a visit to
$[0,H]$ before time $2k-M$ each step moves up with probability at least
$1-1/(4x)$; the visit length is dominated by $H$ plus a geometric
variable whose $x$-th power has finite mean $g_x\le2$ (since
$x\cdot\tfrac1{4x}\le\tfrac14$).  The number of distinct visits is one
plus the number of down-crossings of level $H{+}1$, dominated by
independent indicators with success probabilities $(H{+}1)/(2k-j)$ and
total mean $\le(H{+}1)\log(2k)$; the final $M$ positions contribute at
most $M=O_x(1)$ deterministically.  Hence
$\mathbb E[x^{\Lambda}]\le x^{M+H}
\exp\bigl((H{+}1)\log(2k)(x^{H}g_x-1)\bigr)=k^{O_{x,H}(1)}$, and for
$y\in(0,1]$, $\mathbb E[y^{\Lambda}]\ge
y^{C_H\log k}\,\mathbb P(\Lambda\le C_H\log k)\ge k^{-O_{y,H}(1)}$ by
the same estimates.  Taking $k$-th roots and letting
$\varepsilon\downarrow0$ completes the proof.
\end{proof}

\subsection*{The next free-divisibility target, and why it is delicate}

Equation~(13.2) turns the order-sixteen conditional-positivity test into
a threshold on $m_{16}$.  With the moments now known through order
twenty-two, that test and its successors at orders eighteen, twenty and
twenty-two are all decided in the affirmative: for the odd block indexed
by $\{1,3,5,7,9,11\}$ the Schur slack at order twenty-two is
$\tfrac{2053548476285965691920969192182281957}
{3193518753771609314977865893800}\approx6.43\times10^{5}>0$.  The first
undecided test is the one whose new diagonal entry is $r_{24}$.  Carrying
out the same elimination on the even block indexed by
$\{2,4,6,8,10,12\}$, and using
$r_{24}=m_{24}-\tfrac{12916985465756056}{1091475}$, the condition
$r_{24}\ge v^{\mathsf T}A^{-1}v$ becomes

\begin{equation}\label{eq:m24}
m_{24}(\gamma_T)\;\ge\;
\frac{533385290252268206796914811370716813306691285889}
{36005162406403257088495500297206628000}
\;\approx\;1.4814133713\times10^{10}.
\end{equation}

What makes this the sharp form of the question is how close the target
is.  The successive growth ratios $m_{2k+2}/m_{2k}$ are
\[
2.6667,\ 4.1250,\ 5.5030,\ 6.8557,\ 8.2065,\ 9.5639,\ 10.9292,\
12.3010,\ 13.6771,\ 15.0556,
\]
with second differences $1.3574$, $1.3653$, $1.3718$, $1.3761$,
$1.3785$, so linear extrapolation gives a next ratio near $16.44$ and
\[
m_{24}\;\approx\;1.4821\times10^{10},
\]
which exceeds the threshold in \eqref{eq:m24} by about $0.045$ percent,
roughly one part in $2200$.  No extrapolation from eleven moments can
resolve a margin that small, so the order-twenty-four test is genuinely
undecided by anything short of the exact value; the sign of
$m_{24}-\eqref{eq:m24}$ is a single rational comparison that settles
whether $\gamma_T$ fails free infinite divisibility at this order.

\subsection*{The moment is carried by its own top core class}

Writing $r=k-j$ in the Catalan transform
$m_{2k}=\mathrm{Cat}_k+\sum_{r\ge2}\binom{2k}{k-r}\kappa_r$ turns the
binomial into $\binom{2k}{j}$, so the summand is
$T_j=\binom{2k}{j}\kappa_{k-j}$ and
\[
\frac{T_{j+1}}{T_j}=\frac{2k-j}{j+1}\cdot\frac{\kappa_{k-j-1}}{\kappa_{k-j}} .
\]
Under the scaling $\kappa_r\sim r!\,\gamma^r$ of
Remark~\ref{rmk:kappaasym} this is
$\frac{2k-j}{(j+1)(k-j)\gamma}\to\frac{2}{\gamma(j+1)}$ for $j$ fixed:
the transform is a Poisson-like sum concentrated at the boundary
$r\approx k$.  Summing the geometric-factorial tail gives the following.

\begin{proposition}[Boundary-layer form of the transform]\label{prop:boundary}
If $\kappa_r\sim r!\,\gamma^r$ with $\gamma=2\rho^2$ as conjectured, then
\[
m_{2k}\;\sim\;e^{2/\gamma}\,\kappa_k\;=\;e^{1/\rho^2}\kappa_k ,
\qquad
\frac{m_{2k}}{\mathrm{Cat}_k}\;\sim\;e^{1/\rho^2}\,
\frac{\kappa_k}{\mathrm{Cat}_k}.
\]
In particular the mixing moments of Conjecture~\ref{conj:mix} are the
primitive masses up to the Catalan normalisation and one explicit
constant, so identifying the mixing law is the same problem as
determining the primitive-volume asymptotics.
\end{proposition}

The whole moment is thus carried by the single largest core class: the
empty-core (semicircle) share of $m_{2k}$ falls from $0.231$ at $k=4$ to
$0.004$ at $k=8$, while the share from cores of size $r>k/2$ rises
$0.460$, $0.706$, $0.735$, $0.884$, $0.896$.  The exact ratios
$m_{2k}/\kappa_k$ for $k=2,\dots,8$ are
\[
4.0000,\quad 5.5000,\quad 5.1011,\quad 4.9209,\quad 4.8068,\quad
4.7252,\quad 4.6630,
\]
decreasing monotonically; first-order Richardson extrapolation gives
$4.20$--$4.24$ against the predicted $e^{1/\rho^2}=4.2864$.  This is a
consistency check rather than a measurement: the extrapolation assumes a
$1/k$ correction and uses only $k\le8$.  What it shows is that a
constant obtained independently -- from the Sen--Vir\'ag norm constant
through Conjecture~\ref{conj:slope} -- lands inside the extrapolation's
own spread.

It also explains why finite-order expansions have not identified either
law: every fixed-$r$ term of the transform carries a Beta
representation, hence contributes moments of a law supported in $[0,1]$,
whereas the true mixing law has unbounded support with an exponential
tail.  The unbounded part lives entirely in the boundary layer, which no
fixed-order truncation sees.

\subsection*{A quantitative exclusion for D-finiteness}

Guess-and-verify with held-out moments gives a size-explicit exclusion.
For an ansatz with $u$ unknown coefficients, impose the first $u$
coefficient equations, solve exactly for the nullspace, and test every
solution against the remaining known moments; with $N$ known
coefficients this is a genuine falsification test whenever $u\le N-1$.

Using all twelve known Toeplitz coefficients (through order $22$) and all
ten Hankel ones (through order $18$): every algebraic ansatz
$P(z,M)=0$ of bidegree up to $(6,5)$ and every holonomic ansatz
$\sum_{j\le4}p_j(z)M^{(j)}=0$ of degree up to $6$ that is small enough to
be falsifiable IS falsified -- $13$ algebraic and $12$ holonomic sizes for
$\gamma_T$, $11$ and $10$ for $\gamma_H$, with none consistent with the
known moments.  For instance, at order $14$ the two uniquely determined
ansatzes for $\gamma_H$ predict $10894.88$ and $10866.32$ against the
exact $m_{16}=\tfrac{9179819}{840}=10928.356$.

This is evidence against D-finiteness of either series, not a proof; but
it is consistent with the independent finding that a compressed
recurrence for the volume-weighted series fails a low-order Wick check.
Both series are factorially divergent -- $m_{2k}\asymp(2k-1)!!\,\rho^{2k}$
under Conjecture~\ref{conj:slope} -- so the natural analytic object is
the Borel transform, and the unified conjecture reads: the Borel series
$\sum_k m_{2k}\zeta^k/k!$ and $\sum_r\kappa_r\zeta^r/r!$ have radius of
convergence $1/(2\rho^2)$.
